\setsubtitle{
    author : tobiichi3227\\
    首殺 : 李柏霆
}
\section{G. Performance Analysis}

\begin{frame}{\btitle{題目敘述}}
    給一張有 $N$ 個 basic blocks、$M$ 條有向邊的 CFG
    程式從入口 $1$ 執行到出口 $N$，總共執行 $R$ 次

    \begin{itemize}
        \item 在邊 $e$ 安裝計數器需花費 $C_e$，之後可直接知道流量 $f_e$
        \item 沒裝計數器的邊，只能利用每個點的 flow conservation 推導
        \item 要讓所有邊的流量都能被\textbf{唯一確定}，並最小化安裝成本
    \end{itemize}
\end{frame}

\begin{frame}{\btitle{子任務}}
    \begin{enumerate}
        \item{$M = N - 1$, 且第 $i$ 條邊為 $i \to i + 1$}
        \item{$C_e = 1$}
        \item{$M \le 20$}
        \item{$M = N$ 且無 self-loop 或重複連接}
        \item{每個點雙連通分量皆為單邊分量，或其中每個頂點的分量內度數皆為 $2$}
        \item{無額外限制}
    \end{enumerate}
\end{frame}

\begin{frame}{\btitle{$M = N - 1$, 且第 $i$ 條邊為 $i \to i + 1$}}
    所有邊本來就是一棵鏈，全部都可以不裝，答案是 $0$
\end{frame}

\begin{frame}{\btitle{$C_e=1$}}
    生成樹能省下 $N-1$，所以必須安裝$M-(N-1)=M-N+1$條邊，答案也是 $M-N+1$
\end{frame}

\begin{frame}{\btitle{$M \le 20$}}
    枚舉哪些邊不裝計數器，再用 DSU 檢查是否形成環，取合法集合中省下成本的最大值

    複雜度 $O(2^M M\alpha(N))$
\end{frame}

\begin{frame}{\btitle{$M = N$ 且無 self-loop 或重複連接}}
    這是一張單環圖

    連通無向圖恰有一個環環上至少一條邊必須裝計數器，
    裝成本最小的那條即可，答案是環上最小 $C_e$
\end{frame}

\begin{frame}{\btitle{每個點雙連通分量皆為單邊分量，或其中每個頂點的分量內度數皆為 $2$}}
    這東西就仙人掌圖

    每個點雙連通分量不是單邊就是簡單環
    每個環至少要裝一條邊，而且不同環的選擇互不影響
    對每個環加上其中最小的 $C_e$

    % 這兩個做法都在做同一件事：
    % \begin{center}
    %     每次破壞環時，盡量保留成本大的邊
    % \end{center}
\end{frame}

\begin{frame}{\btitle{無額外限制}}
    普通點滿足
    \[
        \sum_{e\in in(v)}f_e=\sum_{e\in out(v)}f_e,
    \]
    入口與出口則分別多出、少掉 $R$ 的流量
\end{frame}

\begin{frame}{\btitle{無額外限制}}
    令 $F$ 為\textbf{沒有安裝計數器}的邊集合

    如果所有邊都裝計數器，成本是$C_{all}=\sum_{e\in E}C_e$，
    選擇不測量 $F$ 可以省下$\sum_{e\in F}C_e$

    所以問題變成： 在仍能唯一推導所有未知流量的前提下，找出總權重最大的邊集合 $F$
\end{frame}

\begin{frame}{\btitle{無額外限制}}
    先忽略邊的方向，考慮 $F$ 中任意一棵樹，取一個葉節點 $v$

    \begin{center}
    \begin{tikzpicture}[scale=1.05,
        every node/.style={circle,draw,minimum size=0.7cm},
        known/.style={draw=gray!60}, unknown/.style={very thick,blue}]
        \node (v) at (0,0) {$v$};
        \node (u) at (2,0) {$u$};
        \node (a) at (-1.4,0.9) {};
        \node (b) at (-1.4,-0.9) {};
        \draw[unknown,->] (v) -- node[draw=none,above] {?} (u);
        \draw[known,->] (a) -- node[draw=none,above,sloped] {已知} (v);
        \draw[known,->] (v) -- node[draw=none,below,sloped] {已知} (b);
    \end{tikzpicture}
    \end{center}
\end{frame}

\begin{frame}{\btitle{無額外限制}}
    $v$ 只剩一條未知邊，因此用 $v$ 的流量恆等式可算出它

    算完後刪掉這條邊，剩下的仍是森林，反覆拔葉子即可算完所有未知邊

    入口或出口也一樣，只要在恆等式中加入已知的 $R$
\end{frame}

\begin{frame}{\btitle{無額外限制}}
    忽略方向後，假設未知邊形成一個環

    沿著環選擇適當正負號，將各邊流量同時調整 $\varepsilon$：$f_e\longleftarrow f_e\pm\varepsilon$
    每個環上節點的流入與流出變化會互相抵消，
    所有 flow conservation 方程仍成立，計數器讀值也完全不變
\end{frame}

\begin{frame}{\btitle{無額外限制}}
    題目的 $R\ge 2M$ 與可達性保證能讓每條環邊都有正整數流量，
    因此可取 $\varepsilon=1$，得到另一組合法流量，無法唯一決定

    \begin{alertblock}{環也包含特殊情況}
        self-loop 自己就是一個環，兩條連接相同端點的重邊也能形成環
    \end{alertblock}
\end{frame}

\begin{frame}{\btitle{無額外限制}}
    綜合前面所述：
    \[
        \boxed{\text{所有未知流量可唯一確定}\iff F\text{ 忽略方向後是森林}}
    \]

    因此要找最大權重生成森林，最大化省下的成本
\end{frame}

\begin{frame}{\btitle{無額外限制}}
    題目保證每個 basic block 都在某條入口到出口的路徑上，
    所以忽略方向後整張圖連通，又因為 $C_e>0$，最佳森林一定能擴充成生成樹

    \begin{block}{等價問題}
        找忽略方向後的\textbf{最大權重生成樹}
    \end{block}

    時間複雜度 $O(M\log M)$，空間複雜度 $O(N+M)$
\end{frame}

% \begin{frame}{\btitle{Maximum Kruskal（1/2）}}
%     \begin{enumerate}
%         \item 計算 $C_{all}=\sum_e C_e$
%         \item 將所有邊依 $C_e$ 由大到小排序
%         \item 用 DSU 維護忽略方向後的連通塊
%     \end{enumerate}
% \end{frame}

% \begin{frame}{\btitle{Maximum Kruskal（2/2）}}
%     \begin{enumerate}
%         \setcounter{enumi}{3}
%         \item 若邊 $(u,v)$ 的兩端尚未連通，就把它加入最大生成樹，
%         合併兩端並累加 $saved\mathrel{+}=C_e$
%     \end{enumerate}

%     最後
%     \[
%         \boxed{answer=C_{all}-saved}.
%     \]

%     self-loop 的兩端本來就在同一個集合，因此一定會安裝計數器；
%     重邊則最多選一條進入生成樹
% \end{frame}
